% ============================================================
%   §5 Part 4: per-class gene networks
% ============================================================
\section{Per-class gene networks}
\label{sec:part4}

The first three Parts treat each cell as a node in a spatial graph. \partref{4} flips this around: each cell-type class is treated as a population, and the structure I want to recover is the network of conditional dependencies between genes within that population. The tool of choice is the Gaussian Graphical Model (GGM), fit by the Graphical Lasso \cite{friedman2008glasso}. A non-zero entry $\Theta_{ij}$ in the sparse inverse covariance matrix indicates that genes $i$ and $j$ are conditionally dependent given all the other genes; intuitively, an edge in the gene network represents a direct co-expression relationship that survives controlling for everything else. The question of \partref{4} is: do the 8 cell-type classes carry distinct gene networks, and to what extent do those networks share structure?

A second methodological caveat concerns the distributional assumption of the Graphical Lasso. The estimator assumes that the gene expression vectors are multivariate Gaussian, an assumption that spatial-transcriptomics counts violate: log-normalised counts retain zero-inflation and right-skew, especially for genes with cell-type-restricted expression, and the variance grows with the mean rather than being constant. I attenuate this by restricting to a common-gene subset of broadly variable genes (see below), but the violation is real. A principled semiparametric alternative is the nonparanormal SKEPTIC estimator of Liu et al.\ \cite{liu2012nonparanormal}, which replaces Pearson correlations with Kendall's $\tau$ or Spearman's $\rho$ and is consistent under arbitrary continuous marginal transformations. I opted for the standard Graphical Lasso because the conclusions of \partref{4} are qualitative (hub identity, shared-hubs vs.\ rewired-connections) and remain stable across bootstrap resamples; a nonparanormal refit would be a natural validation step but is not pursued in this report.

Before presenting results, a methodological note is necessary. The Visium HD platform resolves expression at near single-cell scale, but at this resolution mRNA from a target cell and its immediate spatial neighbours becomes difficult to disambiguate: a non-trivial fraction of reads assigned to a cell originates from cells within $\sim 1$ to $2$ neighbours of it. This phenomenon, usually called \emph{diffusion contamination} (or molecular bleed-through), means that the expression vector of a cell is partly a reflection of its spatial neighbourhood. Any conditional-dependence network fit on per-cell expression will therefore reflect a mix of cell-intrinsic regulation and the gene programmes of spatially adjacent cells. The strength of this effect varies across cell types: for cells that sit in a homogeneous spatial context the diffusion reinforces the cell-intrinsic signal, while for cells that sit at the interface between two distinct compartments (or for low-abundance cells whose own transcripts are outnumbered by neighbour transcripts) the diffusion may dominate. I refer to the resulting structure as the cell type's \emph{microenvironmental gene programme}, with the understanding that this captures both intrinsic regulation and neighbourhood signal in proportions that vary by class. This is not a methodological flaw of \partref{4} but a property of the data; I revisit its implications explicitly in the Discussion.

\subsection*{The 378 common-gene subset and the choice of \texorpdfstring{$\alpha$}{alpha}}

A GGM on 16{,}497 genes would be both computationally expensive and uninterpretable: with order $10^8$ candidate edges per class and a sample size of $\sim 10^5$ cells, the model would be heavily over-parameterised. Instead, I restrict the analysis to a common-gene subset: the genes among the 1000 HVGs from \partref{3} that have standard deviation above $0.05$ in \emph{every} one of the 8 cell-type classes. This is 378 genes. Restricting to genes that are broadly variable across classes ensures that any class-specific structure cannot trivially be attributed to a gene being silent in some classes and expressed in others; the differences that survive are differences in network topology.

The regularisation parameter $\alpha$ controls sparsity: higher $\alpha$ gives sparser networks. I selected $\alpha = 0.07$ by sweeping over $\{0.05, 0.07, 0.097, 0.136, \ldots, 1.0\}$ on two contrasting classes (SMC, with the most distinct expression pattern; Tumor cells, with the largest expected gene programme) and targeting roughly 1\% edge density. At $\alpha = 0.07$ the eight networks have 494--980 edges out of $\binom{378}{2} = 71{,}253$ possible, corresponding to 0.7\%--1.4\% density. The full sweep showed that the qualitative findings (hub identity, cross-class Jaccards) are stable in the range $\alpha \in [0.05, 0.1]$, so the choice of $0.07$ is not delicate. Numerical caveats: Endothelial cells, Fibroblasts, and T cells did not reach the Graphical Lasso convergence tolerance within 200 iterations at $\alpha = 0.07$; their precision matrices are still well-defined but the dual gap is non-zero. I report results from the best iterate. For each class I subsample 5000 cells from the training pool (patient-stratified) before fitting.

\subsection*{Network topology varies smoothly across classes}

The eight fitted networks all have right-skewed, long-tailed degree distributions (Figure~\ref{fig:part4_degrees}). Most genes have degree below 10; a small number of hub genes have degree above 20 and concentrate the network's connectivity. This is the typical "scale-free flavour" of biological networks, though I do not claim a power-law fit on degree distributions of this size. The class-to-class variation in degree distribution is modest in shape but substantial in scale: Myeloid cells (\classid{3}) and Tumor (\classid{7}) have the longest tails (max degree above 35), T cells (\classid{6}) have the shortest (max degree $\sim 18$), and Normal epithelium (\classid{4}) and T cells (\classid{6}) have the narrowest distributions overall. The numerical edge counts span a factor of 2 across classes (494 to 980), reflecting differences in how much regulatory structure the Graphical Lasso recovers from each class's expression pattern.

\begin{figure}[h]
    \centering
    \includegraphics[width=1.0\linewidth]{part4_degree_distributions.pdf}
    \caption{Degree distributions of the eight per-class GGMs at $\alpha = 0.07$ on the 378 common genes. All distributions are right-skewed with long tails (a small number of hubs concentrate connectivity). T cells (\classid{6}) have the sparsest network ($n_{\text{edges}} = 494$, shortest tail), consistent with their dispersed, infiltrating spatial distribution; Myeloid cells (\classid{3}) have the densest ($n_{\text{edges}} = 980$, longest tail).}
    \label{fig:part4_degrees}
\end{figure}

\subsection*{Hub genes mix cell-intrinsic and neighbourhood signal}

The top-degree hubs of each per-class GGM are biologically interpretable, but a careful reading reveals a more nuanced picture than a clean separation between cell-intrinsic and neighbourhood-derived signal. In several classes the canonical markers of the cell type are present in the top-15 hubs alongside markers of spatially adjacent cell types; in other classes the cell-intrinsic markers are largely diluted out and the neighbourhood signal dominates. I describe each class in turn, with reference to the verified top-15 hubs ranked by degree, and to canonical marker expression patterns separately validated in the expression data.

\textbf{B cells (\classid{0})} hub on DES, ACTG2, MYH11, COL1A2, SMTN, MYLK, CDC20B, SYNM, ACTA2, PPP1R14A, CNN1 --- a smooth-muscle cell signature. The canonical B-cell marker MS4A1 (CD20) does appear at degree 12, but the dominant hubs are an SMC programme. B cells in tumour tissue frequently localise within tertiary lymphoid structures, which are perivascular and surrounded by smooth-muscle vasculature; the GGM picks up the strong SMC programme of this neighbourhood with the B-cell intrinsic signal compressed to a single survivor (MS4A1) in the hub set.

\textbf{Endothelial cells (\classid{1})} are a low-abundance population whose canonical markers (PECAM1, VWF, KDR, CDH5) do not appear in the top-15 hubs at all. The dominant hubs are a mix of inflammatory (CXCL8, S100A9, S100A8, IL1B, G0S2, MMP1, MMP3, MMP12, IL1RN) and vascular SMC (MYH11, DES, MYLK, SYNM) signatures --- the two cell populations that surround endothelial cells in vascularised tumour tissue. The endothelial GGM has the lowest maximum degree of all eight classes (max degree 20 vs.\ degree 38 in B cells), consistent with both the low cell count and the heterogeneous neighbourhood.

\textbf{Fibroblasts (\classid{2})} hub on CXCL8, COL1A2, S100A9, CXCL5, G0S2, S100A8, PLEK, IL1B, PTGS2, MT2A, MMP1, CCL3 --- a mix of canonical fibroblast/CAF markers (COL1A2 at degree 31, the second highest in the network) and an inflammatory/neutrophil signature (CXCL8 at degree 32, S100A8/A9, IL1B, MMP1, CCL3, G0S2). This is one of the cleaner cases of cell-intrinsic and neighbourhood signal coexisting: COL1A2 is a fibroblast marker, the inflammatory genes likely reflect adjacency to inflammatory myeloid cells in inflamed regions of the tumour microenvironment.

\textbf{Myeloid cells (\classid{3})} hub on G0S2, CXCL8, S100A9, CSF3R, S100A8, SOCS3, BCL2A1, FOS, COL1A2, DUSP1, TREM1, PTGS2, IL1B, CXCL5, SELL --- an activated myeloid/granulocyte signature that matches their cell-intrinsic identity directly. CSF3R, BCL2A1, FOS, SOCS3, TREM1, and SELL are all canonical myeloid markers; the inflammatory chemokines (CXCL8, CXCL5) are also produced by myeloid cells themselves. This is the cleanest case where the cell-intrinsic and neighbourhood signal converge rather than compete: myeloid cells cluster spatially with other myeloid populations, so diffusion reinforces the intrinsic programme.

\textbf{Normal epithelium (\classid{4})} has the most heterogeneous hub set: COL1A2 (at degree 21), SLC26A3 (degree 13), MFAP4, DEFA6, IL7R, PTGS2, IGFBP5, G0S2, ADM, LTB, SRGN, ACTA2, FXYD6, TNXB, CXCL8. SLC26A3 and FXYD6 are canonical intestinal-epithelial markers; DEFA6 is a Paneth-cell marker; COL1A2, MFAP4, and TNXB are stromal markers. The hubs reflect both the cell-intrinsic epithelial programme (with intestinal differentiation markers surviving) and the lamina-propria fibroblast adjacency on which intestinal epithelium sits. This is the only class where matrix proteins and canonical epithelial markers appear together in the top-15.

\textbf{SMC (\classid{5})} hubs on S100A9, DES, G0S2, CXCL8, MYH11, S100A8, ACTG2, SELL, EGR1, FOSB, DUSP1, FOS, CXCR2, MT2A, FCAR. The canonical SMC markers DES, MYH11, and ACTG2 appear at top positions (degrees 27, 24, 23 respectively), confirming the cell-intrinsic SMC programme. The strong co-occurrence with inflammatory hubs (CXCL8, S100A8/A9, FOS, FOSB, EGR1, CXCR2) reflects the perivascular position of SMC, which places them at the interface with extravasating myeloid cells. SMC is therefore the converse of B cells (\classid{0}): in B cells the diffusion was from SMC into B cell, here the diffusion is from inflammatory cells into SMC, but the neighbourhood relationship is the same.

\textbf{T cells (\classid{6})} have the most distinctive gene network of all eight classes. Their top hubs are CTSG, MMP25, F13A1, DEFA6, KLHL4, AOC3, IGHM, RNF150, SFRP2, CHRDL2, PNCK, CHIT1, CLCA1, IGLC1, REG3A --- a mixed mast-cell/granulocyte/plasma-cell/Paneth-cell signature. No canonical T-cell marker (CD3D, CD3E, CD8A, CD4, TRAC) appears in the top-15. T cells in tumour tissue are infiltrating populations: they appear scattered in low density throughout multiple compartments, with no consistent spatial neighbourhood. The "mixed" hub signature is a consequence of this spatial heterogeneity --- the GGM averages over many distinct microenvironments and the surviving hubs are those that appear across multiple unusual neighbourhoods, not the common signal of any single one. T cells are also one of the classes for which the Graphical Lasso did not converge in 200 iterations, consistent with heterogeneity-driven instability. Notably the maximum degree in the T-cell network is only 10 (vs.\ 38 in B cells, 41 in Myeloid cells), the lowest of all classes, reflecting the lack of a consistent dominant signal.

\textbf{Tumor (\classid{7})} hubs on CXCL8, COL1A2, IGFBP5, TIMP3, ELN, SFRP4, IL1B, S100A9, MGP, MT2A, SRGN, ACTG2, SCG2, SIGLEC5, CCL5 --- a CAF-rich extracellular matrix signature (COL1A2, TIMP3, ELN, MGP, SFRP4) mixed with inflammatory mediators (CXCL8, IL1B, S100A9, CCL5). Canonical tumour-cell markers (EPCAM, KRT8) are present in the expression data but do not dominate the GGM hubs. Tumor cells in colorectal cancer are typically embedded in CAF-rich stroma with active inflammatory infiltration; both signals coexist in the conditional-dependence network.

A few hubs appear across multiple classes, suggesting genes that are broadly expressed throughout the tumour microenvironment: CXCL8 appears as a hub in 7 of 8 classes, COL1A2 in 6 of 8, S100A8 and S100A9 in 6 of 8, G0S2 in 5 of 8, MYH11 in 4 of 8. These are inflammatory and matrix-remodelling signatures that contaminate many spatial neighbourhoods through diffusion or are genuinely expressed by multiple cell populations; the GGM cannot distinguish the two without an independent measure of cell-intrinsic transcription.
\begin{figure}[p]
    \centering
    \includegraphics[width=\linewidth]{part4_gene_networks.png}
    \caption{Gene networks of two contrasting cell-type classes: Myeloid cells (\classid{3}, left) and B cells (\classid{0}, right). \textbf{Top row}: full GGM on the 378 common genes, drawn as a force-directed network. Edge thickness and opacity are log-scaled by $|\Theta_{ij}|$, so the strongest conditional dependencies emerge from the cloud of weaker edges. Nodes are coloured by gene family (see legend); the top 15 hubs by degree are labelled. \textbf{Bottom row}: the same network reduced to its top 35 hubs by degree, displayed as a chord diagram sorted by family. The Myeloid network is dominated by inflammatory hubs (orange) that reinforce both cell-intrinsic identity and neighbourhood signal; the B-cell network is dominated by SMC markers (red), with MS4A1 (the only canonical B-cell hub) appearing as a single isolated blue node, consistent with diffusion of mRNA from spatially adjacent smooth muscle cells in perivascular tertiary-lymphoid microenvironments.}
    \label{fig:part4_gene_networks}
\end{figure}


Figure~\ref{fig:part4_gene_networks} visualises the contrast between two representative classes: Myeloid cells (\classid{3}) and B cells (\classid{0}). The top row shows each full GGM on all 378 common genes; edge thickness is log-scaled by the magnitude of the inverse covariance entry, so the strongest conditional dependencies emerge as visible threads through the cloud of weaker edges. The bottom row reduces each network to its top 35 hubs and lays them out in a circular order sorted by biological family. The asymmetry between the two classes is striking: Myeloid hubs concentrate in an inflammatory arc (orange) where the cell-intrinsic and neighbourhood signals reinforce each other, while B-cell hubs are dominated by SMC markers (red) with MS4A1, the only canonical B-cell hub, sitting in isolation as the lone blue node next to the SMC cluster. The figure is the cleanest visual rendition of the diffusion-mediated reading of \partref{4}.
\newpage
\subsection*{Networks share neighbourhoods but rewire connections}

The question now is whether the eight networks resemble each other. I use two pairwise similarity measures, computed on the off-diagonals of $8 \times 8$ matrices: the Jaccard overlap of the top-15 hubs by degree (\emph{hub Jaccard}) and the Jaccard overlap of the entire edge set (\emph{edge Jaccard}). Figure~\ref{fig:part4_jaccard} shows both.

\begin{figure}[h]
    \centering
    \includegraphics[width=1.0\linewidth]{part4_jaccard_heatmaps.pdf}
    \caption{Pairwise hub and edge Jaccard between the 8 per-class GGMs. Left: hub Jaccard on top-15 hubs (max off-diagonal $0.36$). Right: edge Jaccard on the full edge sets (max off-diagonal $0.09$). Diagonals are masked. The hub-level overlap is up to four times larger than the edge-level overlap, indicating that cell types share microenvironmental hubs but not the conditional-dependence relationships between them. T cells (\classid{6}) are a structural outlier, with Jaccard $\leq 0.03$ to every other class.}
    \label{fig:part4_jaccard}
\end{figure}

Two observations dominate. First, the maximum hub Jaccard among any pair of classes is $0.36$ (Fibroblasts vs.\ Myeloid cells), and most pairs lie in the range $0.07$--$0.30$. Cell types share microenvironmental hubs in moderate but non-trivial proportions. The cluster of B cells (\classid{0}), Fibroblasts (\classid{2}), Myeloid cells (\classid{3}), and SMC (\classid{5}) is clearly visible as a block of mutual hub overlap: these are the cell types that share the inflammatory/stromal microenvironment of the tumour stroma. Tumor (\classid{7}) and Normal epithelium (\classid{4}) sit somewhat apart from this block; T cells (\classid{6}) are essentially disconnected from all of them.

Second, the maximum edge Jaccard among any pair is only $0.09$. The same Fibroblast--Myeloid pair that shares 36\% of hubs shares only 9\% of edges. Across the full set of pairs, edge Jaccard is roughly an order of magnitude smaller than hub Jaccard. The conclusion is the cleanest finding of \partref{4}:

\begin{insight}
Cell types share hub genes that reflect spatial neighbourhoods (hub overlap up to 36\%) but rewire the specific co-expression relationships between them (edge overlap up to 9\%). In some classes (Myeloid, Fibroblasts, SMC, Tumor, Normal epithelium) the cell-intrinsic markers survive as hubs alongside the neighbourhood signal; in others (B cells, Endothelial, T cells) the neighbourhood dominates. The class-specific signal lives in the edges, not in the nodes.
\end{insight}

This is biologically sensible. Genes like CXCL8 and S100A8/A9 are broadly distributed inflammatory mediators across the tumour microenvironment, and they appear as hubs in any class that sits near an inflammatory neighbourhood. The specific partners they connect to, however, depend on which other genes are co-active in that microenvironment. A B-cell's "inflammatory hub" is conditioned on the SMC-rich vascular context of tertiary lymphoid structures; a myeloid cell's "inflammatory hub" is conditioned on other myeloid neighbours. The hubs are shared; the edges differ. The "rewired connections" finding is the genuine class-specific signal; the "shared regulators" finding is, in this dataset, a partial artefact of diffusion combined with a small set of broadly active inflammatory and matrix-remodelling programmes.

\subsection*{T cells are a microenvironmental outlier}

T cells are the most striking outlier in Figure~\ref{fig:part4_jaccard}: their row and column are essentially 0 on both panels. Hub Jaccard between T cells and the other 7 classes is at most $0.03$ (one shared hub out of 15 with Normal epithelium, the gene DEFA6); edge Jaccard is at most $0.02$. Whatever microenvironmental structure T cells carry, it does not overlap with any of the other classes.

Under the diffusion-contamination framing, this is the consequence of T cells' \emph{lack of a consistent spatial neighbourhood}. T cells in tumour tissue are infiltrating and dispersed: they appear in low density spread across every compartment of the tissue, with no consistent partner cell type. The result is that their microenvironmental gene programme averages over many distinct local contexts, and the dominant signal that survives is the one shared across all the unusual neighbourhoods that T cells happen to find themselves in --- not the common signature of any one neighbourhood. The convergence failure of the Graphical Lasso on this class is consistent with this picture: heterogeneous neighbourhoods produce unstable covariance estimates.

This re-interpretation of the T-cell outlier --- microenvironmental dispersion rather than cell-intrinsic distinctness --- aligns with the topological role identified in \partref{1} (T cells sit at the boundary between dense and sparse regions, with negative degree assortativity) and the niche role identified in \partref{2} (T cells never dominate a niche but appear as a secondary minority in \nicheid{17} only). All three signatures point to the same underlying biology: T cells are a spatially diffuse population, and any per-class analysis that depends on consistent spatial context will flag them as distinctive.

\subsection*{Bootstrap stability}

To confirm that the headline findings are not an artefact of the particular 5000-cell subsample, I refit each class network on 10 bootstrap resamples (sampled with replacement from the training pool) and tracked hub identity, edge counts, and edge stability. The canonical hub genes (DES, MYH11, ACTA2 for B cells; CXCL8, S100A8/9 for Fibroblasts; COL1A2, TIMP3 for Tumor) have coefficient of variation below $0.20$ across the 10 bootstraps. Edge counts reproduce within $\pm 3$--$5\%$ of the headline value. Edge identity is less stable: only 22\%--34\% of the headline edges appear in at least 8 of 10 bootstraps. The conclusion is that the microenvironmental hubs are stable, edge counts are stable, but individual edges have non-trivial variance --- as expected from a sparse-precision-matrix estimator at moderate sample size.

\medskip

\noindent\textbf{Summary.} Per-class GGMs at $\alpha = 0.07$ on a 378-gene common subset recover gene networks whose top-degree hubs mix cell-intrinsic markers and diffusion-mediated signal from spatially adjacent cell types. Some classes (Myeloid, Fibroblasts, SMC, Tumor, Normal epithelium) retain canonical cell-type markers among their hubs; others (B cells, Endothelial, T cells) are dominated by neighbourhood signal. Across the 8 networks, hub Jaccard reaches $0.36$ (Fibroblasts vs.\ Myeloid cells) while edge Jaccard stays below $0.09$: classes share hubs but rewire connections. T cells are a structural outlier (hub and edge Jaccard $\leq 0.03$ to all others), consistent with their dispersed infiltrating spatial distribution and the absence of a consistent neighbourhood.